\ExerciseSection

\begin{enumerate}
\def\labelenumi{\arabic{enumi})}
\tightlist
\item
  A pseudometric \(f\) on a group \(G\) (written multiplicatively) is said to be left-invariant (resp. right-invariant) if \(f(zx, zy) = f(x, y)\) {[}resp. \(f(xz, yz) = f(x, y)\) {]} for all \(x, y, z\) in \(G\) .
\end{enumerate}

\begin{enumerate}
\def\labelenumi{\alph{enumi})}
\tightlist
\item
  If \(f\) is a left-invariant pseudometric, the real-valued function \(g(x) = f(e, x)\) on \(\mathbf{G}\) ( \(e\) being the identity element of \(\mathbf{G}\) ) satisfies the following conditions:
\end{enumerate}

\begin{enumerate}
\def\labelenumi{(\roman{enumi})}
\item
  \(g(x)\geqslant \mathrm{0}\) for all \(x\in \mathbf{G}\) , and \(g(e) = \mathrm{0};\)
\item
  \(g(x^{-1}) = g(x)\) ;
\item
  \(g(xy) \leqslant g(x) + g(y)\) .
\end{enumerate}

Conversely, if \(g\) is any real-valued function on \(G\) satisfying these conditions, then \(f(x, y) = g(x^{-1}y)\) is a left-invariant pseudometric on \(G\) . \(b)\) The topology \(\mathfrak{G}\) defined by a saturated family \((f_t)\) of left-invariant pseudometrics on a group \(G\) is compatible with the group structure of \(G\) if and only if, for each \(a \in G\) , each index \(t\) and each real number \(\alpha > 0\) , there exists an index \(x\) and a real number \(\beta > 0\) such that the relation \(f_x(e, x) \leqslant \beta\) implies \(f_t(e, axa^{-1}) \leqslant \alpha\) . If this condition is satisfied, the uniformity on \(G\) defined by the family \(f_t\) is the same as the left uniformity of the topological group obtained by endowing \(G\) with the topology \(\mathfrak{G}\) .

\begin{enumerate}
\def\labelenumi{\alph{enumi})}
\setcounter{enumi}{2}
\tightlist
\item
  On every topological group G there exists a family of left-invariant pseudometrics such that the uniformity defined by this family is the same as the left uniformity of G.
\end{enumerate}

\begin{enumerate}
\def\labelenumi{\arabic{enumi})}
\setcounter{enumi}{1}
\item
  Let \(\mathbf{G}\) be a topological group whose left and right uniformities coincide. Show that this unique uniformity can be defined by a family of pseudometrics which are simultaneously left-and right-invariant (using Exercise 3 of Chapter III, § 3, show that if \(\mathbf{V}\) is any neighbourhood of the identity element \(e\) of \(\mathbf{G}\) , then \(\mathbf{V}_0 = \bigcap_{x\in \mathbf{G}}x\mathbf{V}x^{-1}\) is a neighbourhood of \(e\) ).
\item
  Let G be a topological group, let \((f_{t})\) be a saturated family of left-invariant pseudometrics on G which define the left uniformity of G, and let \(g_{t}(x)=f_{t}(e,x)\) . Let H be a normal closed subgroup of G, and for each coset \(\dot{x}\in G/H\) let \(h_{t}(\dot{x})=\inf_{x\in \dot{x}}g_{t}(x)\) ; show that, if \(\bar{f}_{t}(\dot{x},\dot{y})=h_{t}(\dot{x}^{-1}\dot{y})\) , the \(\bar{f}_{t}\) form a family of left-invariant pseudometrics on G/H which define the left uniformity of G/H (argue as in Remark 2 of no. 1).
\item
  Show that the topology of a valued division ring is locally retrobounded (Chapter III, § 6, Exercise 22).
\end{enumerate}

¶ 5) Let \(\omega\) be an absolute value on the field \(\mathbf{Q}\) of rational numbers. a) Show that \(\omega\) is uniquely determined by its values at the prime numbers. b) If there is a prime number \(p\) such that \(\omega(p) \leqslant 1\) , show that \(\omega(q) \leqslant 1\) for every other prime number \(q\) {[}find an upper bound for \(\omega(q^n)\) by writing \(q^n\) in the scale of \(p\) , and let \(n \to \infty\) {]}.

\begin{enumerate}
\def\labelenumi{\alph{enumi})}
\setcounter{enumi}{2}
\tightlist
\item
  If there exists a prime number p such that \(\omega(p) < 1\) , show that \(\omega(q) = 1\) for every other prime number q {[}show that we cannot have \(\omega(q) < 1\) , by using the fact that for each integer n \textgreater{} 0 there exist two rational integers r and s such that \(1 = rp^{n} + sq^{n}\) {]} . Hence show that \(\omega\) is then an absolute value equivalent to the p-adic absolute value on Q. d) If \(\omega(p) > 1\) for every prime number p, show that if p and q are any two prime numbers, then
\end{enumerate}

\[
\frac {\log \omega (p)}{\log p} = \frac {\log \omega (q)}{\log q}
\]

{[}same method as in \(b\) {]}. Hence show that in this case \(\omega(x) = |x|^{\rho}\) for some \(\rho \leqslant 1\) .

\begin{enumerate}
\def\labelenumi{\arabic{enumi})}
\setcounter{enumi}{5}
\item
  Let E be a left vector space over a division ring K, with a countable base \((a_{n})\) . Let \((r_{n})\) be a decreasing sequence of numbers \textgreater0, tending to 0. For each \(x = \sum_{k} t_{k} a_{k} \neq 0\) of E, put \(\|x\| = r_{h}\) where h is the smallest of the indices k such that \(t_{k} \neq 0\) , and put \(\|0\| = 0\) . Show that \(\|x - y\|\) is an invariant metric on the additive group E, and that the topology it defines on E is independent of the sequence \((r_{n})\) (decreasing and tending to 0) chosen. Hence show that if the definitions of a norm and a normed space (Definitions 5 and 6) are extended to the case where the absolute value on the scalars is improper, Proposition 7 and Theorem 1 are no longer valid.
\item
  Let E be a normed space over a valued division ring. Show that if every absolutely summable family in E is summable in E, then E is complete. {[}Let \((x_{n})\) be a Cauchy sequence in E, and consider a subsequence \((x_{n_{k}})\) of \((x_{n})\) such that the series whose general term is \(x_{n_{k+1}} - x_{n}\) is absolutely convergent.{]}
\item
  Give an example of a family which is summable but not absolutely summable in the field \(Q_{p}\) of p-adic numbers. {[}Show that a family \((x_{t})_{t\in\mathbf{I}}\) of points of \(Q_{p}\) is summable if and only if \(\lim x_{t}=0\) with respect to the filter of complements of finite subsets of I.{]}
\item
  A mapping \(w\) of a ring \(\mathbf{A}\) into \(\mathbf{R}_+\) is called a semi-absolute value on \(\mathbf{A}\) if it satisfies the conditions: (i) \(w(\mathrm{0}) = \mathrm{0}\) ; (ii) \(w(x - y) \leqslant w(x) + w(y)\) ;
\end{enumerate}

\begin{enumerate}
\def\labelenumi{(\roman{enumi})}
\setcounter{enumi}{2}
\tightlist
\item
  \(w(xy) \leqslant w(x)w(y)\) . The semi-absolute value w is said to be Hausdorff if \(w(x) = 0\) implies x = 0. If w is a Hausdorff semi-absolute value on A, then \(w(x - y)\) is an invariant metric on the additive group of A, and hence defines a topology compatible with the additive group structure of A. Generalize the principal properties of normed algebras, notably Proposition 13, to rings endowed with a semi-absolute value.
\end{enumerate}

If \(w\) is a non-Hausdorff semi-absolute value on \(A\) , the set \(\overline{w}(0)\) is a two-sided ideal \(\mathfrak{a}\) in \(A\) ; if we endow \(A\) with the topology defined by the pseudometric \(w(x - y)\) , this topology is compatible with the ring structure of \(A\) , and the Hausdorff space associated with \(A\) is the quotient ring \(A / \mathfrak{a}\) , on which the function \(\overline{w}(\overline{x})\) , which for each \(\overline{x} \mod \mathfrak{a}\) is equal to the common value of \(w(x)\) for all \(x \in \overline{x}\) , is a Hausdorff semi-absolute value said to be associated with \(w\) ; the topology defined by \(\overline{w}\) is the quotient by \(\mathfrak{a}\) of the topology of \(A\) .

¶ 10) a) Two semi-absolute values \(w_1, w_2\) on a ring A are said to be equivalent if the pseudometrics \(w_1(x - y), w_2(x - y)\) are equivalent. Show that, if \(w\) is a semi-absolute value on A, then \(aw\) and \(w^{1/a}\) are semi-absolute values equivalent to \(w\) for every real number \(a \geqslant 1\) .

\begin{enumerate}
\def\labelenumi{\alph{enumi})}
\setcounter{enumi}{1}
\tightlist
\item
  If \(w_{i}\) ( \(1 \leqslant i \leqslant n\) ) are semi-absolute values on a ring A, the functions \(w = \sum_{i} w_{i}\) and \(w' = \sup_{i} w_{i}\) are two equivalent semi-absolute values on A. If
\end{enumerate}

\[
\mathfrak {a} _ {i} = \overline {{{w}}} _ {i} ^ {1} (\mathrm{0}) \quad \text { and } \quad \mathfrak {a} = \overline {{{w}}} ^ {1} (\mathrm{0}),
\]

then \(\mathfrak{a} = \bigcap_{i} \mathfrak{a}_{i}\) . If \(A_i\) denotes the completion of the quotient ring \(A/a_i\) endowed with the Hausdorff semi-absolute value associated with \(w_i\) , show that \(A/a_i\) endowed with the Hausdorff semi-absolute value associated with \(w\) is isomorphic to a subring of the product ring \(\prod_{i} A_i\) ; show that (assuming that \(A\) has an identity element \(1\) ) this ring is dense in \(\prod_{i} A_i\) if and only if the \(w_i\) satisfy the following condition: for each \(\varepsilon > 0\) and for each index \(i\) , there is an element \(x_i \in A\) such that \(w_i(1 - x_i) \leqslant \varepsilon\) and \(w_k(x_i) \leqslant \varepsilon\) for each index \(k \neq i\) .

\begin{enumerate}
\def\labelenumi{\arabic{enumi})}
\setcounter{enumi}{10}
\tightlist
\item
  Let A be a ring endowed with a Hausdorff semi-absolute value, and suppose that A has an identity element and is complete with respect to the topology defined by the semi-absolute value. Show that every maximal ideal of A is closed (use Proposition 13).
\end{enumerate}

¶ 12) Let K be a non-discrete, Hausdorff, topological division ring, and let \(\varphi: \mathbf{K} \to \mathbf{R}_{+}\) be a mapping such that \(\varphi(0) = 0\) , \(\varphi(xy) = \varphi(x)\varphi(y)\) and such that if \(V_{n}\) denotes the set of all \(x \in K\) such that \(\varphi(x) \leqslant 1/n\) , the sets \(V_{n}\) form a fundamental system of neighbourhoods of \(0\) in K.

\begin{enumerate}
\def\labelenumi{\alph{enumi})}
\tightlist
\item
  Show that there exists a real number \(a > 0\) such that, for all \(x \in \mathbf{K}\) ,
\end{enumerate}

\[
\varphi (1 + x) \leqslant a (1 + \varphi (x))
\]

{[}if not, there would exist a sequence \((x_{n})\) of points of K such that both \((1 + x_n)^{-1}\) and \(x_{n}(1 + x_{n})^{-1}\) tended to zero{]}. Hence show that if \(\psi (x) = (\varphi (x))^{\alpha}\) we have

\[
\psi (x + y) \leqslant 2 \sup (\psi (x), \psi (y))
\]

for \(\alpha\) sufficiently small.

b) Show that if \(n = 2^p\) we have \(\psi\left(\sum_{i=1}^{n} x_i\right) \leqslant n \sup (\psi(x_i))\) , and deduce that for every integer \(m > 0\) we have \(\psi(m) \leqslant 2m\) .

\begin{enumerate}
\def\labelenumi{\alph{enumi})}
\setcounter{enumi}{2}
\tightlist
\item
  Deduce from b) that for every \(n = 2^p\) and every \(x \in \mathbf{K}\) we have \(\psi((1 + x)^{n-1}) \leqslant 2n(1 + \psi(x))^{n-1}\) , and hence that \(\psi\) is an absolute value on \(\mathbf{K}\) which defines the topology of \(\mathbf{K}\) .
\end{enumerate}

\begin{enumerate}
\def\labelenumi{\arabic{enumi})}
\setcounter{enumi}{12}
\tightlist
\item
  Let K be a non-discrete Hausdorff topological division ring. Let R be the set of all \(x \in K\) such that \(\lim_{n \to \infty} x^n = 0\) , and let N be the complement of the set \(R \cup R^{-1}\) . Show that there exists an absolute value on K which defines the topology of K if and only if (i) R is open in K; (ii) for each neighbourhood V of 0 in K, there is a neighbourhood U of 0 in K such that \(R.U \subset V\) ; and (iii) if \(x \in R\) and \(y \in R \cup N\) , then \(yx \in R\) .
\end{enumerate}

To show that these conditions are sufficient, prove successively that: a) N is a normal subgroup of the multiplicative group \(K^{*}\) of non-zero elements of K.

\begin{enumerate}
\def\labelenumi{\alph{enumi})}
\setcounter{enumi}{1}
\item
  In the quotient group \(K^{*}/N\) , put \(\dot{x} \leqslant \dot{y}\) if there exist \(x \in \dot{x}\) and \(y \in \dot{y}\) such that \(yx^{-1} \in R \cup N\) ; show that this relation is an ordering compatible with the group structure of \(K^{*}/N\) , and that the ordered group so defined is isomorphic to a subgroup of the additive group R (use Exercise 1 of Chapter V, § 3).
\item
  Complete the proof with the help of Exercise 12.
\end{enumerate}

In particular, the topology of a non-discrete locally compact topological field can be defined by an absolute value.

